\documentclass[11pt]{article}
\newcommand{\videoid}{6GBBu7gR3tY}
% Preámbulo compartido de las notas del curso de topología algebraica.
% Cada nota empieza con \documentclass[11pt]{article}, define \videoid{...}
% (el identificador del video de YouTube) y luego hace % Preámbulo compartido de las notas del curso de topología algebraica.
% Cada nota empieza con \documentclass[11pt]{article}, define \videoid{...}
% (el identificador del video de YouTube) y luego hace % Preámbulo compartido de las notas del curso de topología algebraica.
% Cada nota empieza con \documentclass[11pt]{article}, define \videoid{...}
% (el identificador del video de YouTube) y luego hace \input{preambulo}.
% Compilar con LuaLaTeX:  lualatex nombre-de-la-nota.tex
\usepackage[spanish,provide=*]{babel}
\usepackage[T1]{fontenc}
\usepackage{lmodern}
\usepackage[a4paper,margin=2.6cm]{geometry}
\usepackage{amsmath,amssymb,amsthm}
\usepackage{xcolor}
\usepackage{tikz}
\usetikzlibrary{arrows.meta,babel}
\usepackage[colorlinks=true,urlcolor=blue!50!black,linkcolor=blue!50!black]{hyperref}

\hyphenpenalty=10000 \exhyphenpenalty=10000 \emergencystretch=3em

% Colores de las figuras (los mismos que usa la página)
\definecolor{extremo}{RGB}{47,107,79}
\definecolor{azul}{RGB}{44,74,154}

% Definiciones numeradas por unidad
\theoremstyle{definition}
\newtheorem{defn}{Definición}
% Numeración por unidad: \unidad{u}{n} (después de \input{preambulo}) hace que las
% definiciones de la nota se llamen u.n, u.(n+1), ...; n es el primer número
% libre de esa unidad (ver la tabla de CONVENCIONES.md).
\newcommand{\unidad}[2]{\renewcommand{\thedefn}{#1.\arabic{defn}}\setcounter{defn}{\numexpr#2-1\relax}}

% Observaciones y erratas del video. Usar SOLO si hay algo que señalar:
%   \begin{erratas}
%     \errata{203}{3:23}{Se dice ..., pero debe decir ...}
%   \end{erratas}
% El primer argumento son los segundos (enlaza al minuto exacto del video),
% el segundo es la marca visible, el tercero es el texto.
\newenvironment{erratas}{\section*{Observaciones y erratas del video}\begin{itemize}}{\end{itemize}}
\newcommand{\errata}[3]{%
  \edef\errurl{https://www.youtube.com/watch?v=\videoid\string&t=#1}%
  \item[\textbf{\expandafter\href\expandafter{\errurl}{#2}}] #3}
.
% Compilar con LuaLaTeX:  lualatex nombre-de-la-nota.tex
\usepackage[spanish,provide=*]{babel}
\usepackage[T1]{fontenc}
\usepackage{lmodern}
\usepackage[a4paper,margin=2.6cm]{geometry}
\usepackage{amsmath,amssymb,amsthm}
\usepackage{xcolor}
\usepackage{tikz}
\usetikzlibrary{arrows.meta,babel}
\usepackage[colorlinks=true,urlcolor=blue!50!black,linkcolor=blue!50!black]{hyperref}

\hyphenpenalty=10000 \exhyphenpenalty=10000 \emergencystretch=3em

% Colores de las figuras (los mismos que usa la página)
\definecolor{extremo}{RGB}{47,107,79}
\definecolor{azul}{RGB}{44,74,154}

% Definiciones numeradas por unidad
\theoremstyle{definition}
\newtheorem{defn}{Definición}
% Numeración por unidad: \unidad{u}{n} (después de % Preámbulo compartido de las notas del curso de topología algebraica.
% Cada nota empieza con \documentclass[11pt]{article}, define \videoid{...}
% (el identificador del video de YouTube) y luego hace \input{preambulo}.
% Compilar con LuaLaTeX:  lualatex nombre-de-la-nota.tex
\usepackage[spanish,provide=*]{babel}
\usepackage[T1]{fontenc}
\usepackage{lmodern}
\usepackage[a4paper,margin=2.6cm]{geometry}
\usepackage{amsmath,amssymb,amsthm}
\usepackage{xcolor}
\usepackage{tikz}
\usetikzlibrary{arrows.meta,babel}
\usepackage[colorlinks=true,urlcolor=blue!50!black,linkcolor=blue!50!black]{hyperref}

\hyphenpenalty=10000 \exhyphenpenalty=10000 \emergencystretch=3em

% Colores de las figuras (los mismos que usa la página)
\definecolor{extremo}{RGB}{47,107,79}
\definecolor{azul}{RGB}{44,74,154}

% Definiciones numeradas por unidad
\theoremstyle{definition}
\newtheorem{defn}{Definición}
% Numeración por unidad: \unidad{u}{n} (después de \input{preambulo}) hace que las
% definiciones de la nota se llamen u.n, u.(n+1), ...; n es el primer número
% libre de esa unidad (ver la tabla de CONVENCIONES.md).
\newcommand{\unidad}[2]{\renewcommand{\thedefn}{#1.\arabic{defn}}\setcounter{defn}{\numexpr#2-1\relax}}

% Observaciones y erratas del video. Usar SOLO si hay algo que señalar:
%   \begin{erratas}
%     \errata{203}{3:23}{Se dice ..., pero debe decir ...}
%   \end{erratas}
% El primer argumento son los segundos (enlaza al minuto exacto del video),
% el segundo es la marca visible, el tercero es el texto.
\newenvironment{erratas}{\section*{Observaciones y erratas del video}\begin{itemize}}{\end{itemize}}
\newcommand{\errata}[3]{%
  \edef\errurl{https://www.youtube.com/watch?v=\videoid\string&t=#1}%
  \item[\textbf{\expandafter\href\expandafter{\errurl}{#2}}] #3}
) hace que las
% definiciones de la nota se llamen u.n, u.(n+1), ...; n es el primer número
% libre de esa unidad (ver la tabla de CONVENCIONES.md).
\newcommand{\unidad}[2]{\renewcommand{\thedefn}{#1.\arabic{defn}}\setcounter{defn}{\numexpr#2-1\relax}}

% Observaciones y erratas del video. Usar SOLO si hay algo que señalar:
%   \begin{erratas}
%     \errata{203}{3:23}{Se dice ..., pero debe decir ...}
%   \end{erratas}
% El primer argumento son los segundos (enlaza al minuto exacto del video),
% el segundo es la marca visible, el tercero es el texto.
\newenvironment{erratas}{\section*{Observaciones y erratas del video}\begin{itemize}}{\end{itemize}}
\newcommand{\errata}[3]{%
  \edef\errurl{https://www.youtube.com/watch?v=\videoid\string&t=#1}%
  \item[\textbf{\expandafter\href\expandafter{\errurl}{#2}}] #3}
.
% Compilar con LuaLaTeX:  lualatex nombre-de-la-nota.tex
\usepackage[spanish,provide=*]{babel}
\usepackage[T1]{fontenc}
\usepackage{lmodern}
\usepackage[a4paper,margin=2.6cm]{geometry}
\usepackage{amsmath,amssymb,amsthm}
\usepackage{xcolor}
\usepackage{tikz}
\usetikzlibrary{arrows.meta,babel}
\usepackage[colorlinks=true,urlcolor=blue!50!black,linkcolor=blue!50!black]{hyperref}

\hyphenpenalty=10000 \exhyphenpenalty=10000 \emergencystretch=3em

% Colores de las figuras (los mismos que usa la página)
\definecolor{extremo}{RGB}{47,107,79}
\definecolor{azul}{RGB}{44,74,154}

% Definiciones numeradas por unidad
\theoremstyle{definition}
\newtheorem{defn}{Definición}
% Numeración por unidad: \unidad{u}{n} (después de % Preámbulo compartido de las notas del curso de topología algebraica.
% Cada nota empieza con \documentclass[11pt]{article}, define \videoid{...}
% (el identificador del video de YouTube) y luego hace % Preámbulo compartido de las notas del curso de topología algebraica.
% Cada nota empieza con \documentclass[11pt]{article}, define \videoid{...}
% (el identificador del video de YouTube) y luego hace \input{preambulo}.
% Compilar con LuaLaTeX:  lualatex nombre-de-la-nota.tex
\usepackage[spanish,provide=*]{babel}
\usepackage[T1]{fontenc}
\usepackage{lmodern}
\usepackage[a4paper,margin=2.6cm]{geometry}
\usepackage{amsmath,amssymb,amsthm}
\usepackage{xcolor}
\usepackage{tikz}
\usetikzlibrary{arrows.meta,babel}
\usepackage[colorlinks=true,urlcolor=blue!50!black,linkcolor=blue!50!black]{hyperref}

\hyphenpenalty=10000 \exhyphenpenalty=10000 \emergencystretch=3em

% Colores de las figuras (los mismos que usa la página)
\definecolor{extremo}{RGB}{47,107,79}
\definecolor{azul}{RGB}{44,74,154}

% Definiciones numeradas por unidad
\theoremstyle{definition}
\newtheorem{defn}{Definición}
% Numeración por unidad: \unidad{u}{n} (después de \input{preambulo}) hace que las
% definiciones de la nota se llamen u.n, u.(n+1), ...; n es el primer número
% libre de esa unidad (ver la tabla de CONVENCIONES.md).
\newcommand{\unidad}[2]{\renewcommand{\thedefn}{#1.\arabic{defn}}\setcounter{defn}{\numexpr#2-1\relax}}

% Observaciones y erratas del video. Usar SOLO si hay algo que señalar:
%   \begin{erratas}
%     \errata{203}{3:23}{Se dice ..., pero debe decir ...}
%   \end{erratas}
% El primer argumento son los segundos (enlaza al minuto exacto del video),
% el segundo es la marca visible, el tercero es el texto.
\newenvironment{erratas}{\section*{Observaciones y erratas del video}\begin{itemize}}{\end{itemize}}
\newcommand{\errata}[3]{%
  \edef\errurl{https://www.youtube.com/watch?v=\videoid\string&t=#1}%
  \item[\textbf{\expandafter\href\expandafter{\errurl}{#2}}] #3}
.
% Compilar con LuaLaTeX:  lualatex nombre-de-la-nota.tex
\usepackage[spanish,provide=*]{babel}
\usepackage[T1]{fontenc}
\usepackage{lmodern}
\usepackage[a4paper,margin=2.6cm]{geometry}
\usepackage{amsmath,amssymb,amsthm}
\usepackage{xcolor}
\usepackage{tikz}
\usetikzlibrary{arrows.meta,babel}
\usepackage[colorlinks=true,urlcolor=blue!50!black,linkcolor=blue!50!black]{hyperref}

\hyphenpenalty=10000 \exhyphenpenalty=10000 \emergencystretch=3em

% Colores de las figuras (los mismos que usa la página)
\definecolor{extremo}{RGB}{47,107,79}
\definecolor{azul}{RGB}{44,74,154}

% Definiciones numeradas por unidad
\theoremstyle{definition}
\newtheorem{defn}{Definición}
% Numeración por unidad: \unidad{u}{n} (después de % Preámbulo compartido de las notas del curso de topología algebraica.
% Cada nota empieza con \documentclass[11pt]{article}, define \videoid{...}
% (el identificador del video de YouTube) y luego hace \input{preambulo}.
% Compilar con LuaLaTeX:  lualatex nombre-de-la-nota.tex
\usepackage[spanish,provide=*]{babel}
\usepackage[T1]{fontenc}
\usepackage{lmodern}
\usepackage[a4paper,margin=2.6cm]{geometry}
\usepackage{amsmath,amssymb,amsthm}
\usepackage{xcolor}
\usepackage{tikz}
\usetikzlibrary{arrows.meta,babel}
\usepackage[colorlinks=true,urlcolor=blue!50!black,linkcolor=blue!50!black]{hyperref}

\hyphenpenalty=10000 \exhyphenpenalty=10000 \emergencystretch=3em

% Colores de las figuras (los mismos que usa la página)
\definecolor{extremo}{RGB}{47,107,79}
\definecolor{azul}{RGB}{44,74,154}

% Definiciones numeradas por unidad
\theoremstyle{definition}
\newtheorem{defn}{Definición}
% Numeración por unidad: \unidad{u}{n} (después de \input{preambulo}) hace que las
% definiciones de la nota se llamen u.n, u.(n+1), ...; n es el primer número
% libre de esa unidad (ver la tabla de CONVENCIONES.md).
\newcommand{\unidad}[2]{\renewcommand{\thedefn}{#1.\arabic{defn}}\setcounter{defn}{\numexpr#2-1\relax}}

% Observaciones y erratas del video. Usar SOLO si hay algo que señalar:
%   \begin{erratas}
%     \errata{203}{3:23}{Se dice ..., pero debe decir ...}
%   \end{erratas}
% El primer argumento son los segundos (enlaza al minuto exacto del video),
% el segundo es la marca visible, el tercero es el texto.
\newenvironment{erratas}{\section*{Observaciones y erratas del video}\begin{itemize}}{\end{itemize}}
\newcommand{\errata}[3]{%
  \edef\errurl{https://www.youtube.com/watch?v=\videoid\string&t=#1}%
  \item[\textbf{\expandafter\href\expandafter{\errurl}{#2}}] #3}
) hace que las
% definiciones de la nota se llamen u.n, u.(n+1), ...; n es el primer número
% libre de esa unidad (ver la tabla de CONVENCIONES.md).
\newcommand{\unidad}[2]{\renewcommand{\thedefn}{#1.\arabic{defn}}\setcounter{defn}{\numexpr#2-1\relax}}

% Observaciones y erratas del video. Usar SOLO si hay algo que señalar:
%   \begin{erratas}
%     \errata{203}{3:23}{Se dice ..., pero debe decir ...}
%   \end{erratas}
% El primer argumento son los segundos (enlaza al minuto exacto del video),
% el segundo es la marca visible, el tercero es el texto.
\newenvironment{erratas}{\section*{Observaciones y erratas del video}\begin{itemize}}{\end{itemize}}
\newcommand{\errata}[3]{%
  \edef\errurl{https://www.youtube.com/watch?v=\videoid\string&t=#1}%
  \item[\textbf{\expandafter\href\expandafter{\errurl}{#2}}] #3}
) hace que las
% definiciones de la nota se llamen u.n, u.(n+1), ...; n es el primer número
% libre de esa unidad (ver la tabla de CONVENCIONES.md).
\newcommand{\unidad}[2]{\renewcommand{\thedefn}{#1.\arabic{defn}}\setcounter{defn}{\numexpr#2-1\relax}}

% Observaciones y erratas del video. Usar SOLO si hay algo que señalar:
%   \begin{erratas}
%     \errata{203}{3:23}{Se dice ..., pero debe decir ...}
%   \end{erratas}
% El primer argumento son los segundos (enlaza al minuto exacto del video),
% el segundo es la marca visible, el tercero es el texto.
\newenvironment{erratas}{\section*{Observaciones y erratas del video}\begin{itemize}}{\end{itemize}}
\newcommand{\errata}[3]{%
  \edef\errurl{https://www.youtube.com/watch?v=\videoid\string&t=#1}%
  \item[\textbf{\expandafter\href\expandafter{\errurl}{#2}}] #3}

\unidad{3}{13} % unidad 3; primer enunciado: 3.13

\title{Construcción del cubriente universal}
\author{Luis Jorge Sánchez Saldaña\\[2pt]\small Facultad de Ciencias, UNAM}
\date{}

\begin{document}
\maketitle

\noindent\small Notas del video
\href{https://www.youtube.com/watch?v=6GBBu7gR3tY}{\guillemotleft{}El cubriente universal - parte 1\guillemotright{}}
(curso de Introducción a la Topología Algebraica). Se define el cubriente universal y se encuentra una condición necesaria para que un espacio lo tenga: ser semilocalmente simplemente conexo.
\normalsize

\section*{El programa}

En este video se comienza a hablar del cubriente universal de un espacio. El objetivo final de estos videos es clasificar todos los cubrientes que tiene un espacio; pero no se va a trabajar con un espacio topológico cualquiera, sino que se va a descubrir que se necesitan algunas hipótesis sobre los espacios.

Para comenzar el programa de clasificar cubrientes, una pregunta natural es la siguiente: dado un espacio $X$, ?`existe algún cubriente $p\colon\tilde X\to X$ de manera que $\pi_1(\tilde X)$ sea trivial? Es decir, ?`tiene $X$ un cubriente simplemente conexo? Para que esta pregunta tenga una respuesta afirmativa, primero se hace un análisis preliminar y se descubren condiciones necesarias sobre el espacio $X$ para que tenga un cubriente de este estilo.

\begin{defn}[Cubriente universal]\label{def:cubriente-universal}
Un cubriente $p\colon\tilde X\to X$ con $\tilde X$ simplemente conexo se conoce como el \emph{cubriente universal} de $X$.
\end{defn}

El nombre se justifica por resultados que se van a obtener más adelante, en los que se verá que, esencialmente, este cubriente es universal en el sentido de que cubre a cualquier otro cubriente de $X$.

\medskip\noindent\textbf{Hipótesis sobre $X$.} Para comenzar se supone que $X$ es arco-conexo. Como se comentó en el video en el que se introdujo la definición de cubriente, esencialmente siempre basta con restringirse a espacios conexos; aquí se supone algo más, que $X$ es arco-conexo, porque solamente en este caso los cubrientes son interesantes, y a la hora de clasificar se van a querer clasificar cubrientes que también son arco-conexos. Por razones un poco más técnicas, que se verán después, también se supone que $X$ es localmente arco-conexo (Definición~3.10): si se tiene un punto $x\in X$ y una vecindad $U$ de $x$, siempre existe otra vecindad más pequeña $V$, con $x\in V\subseteq U$, que es arco-conexa.

\section*{Una condición necesaria}

Supóngase que ya existe un cubriente $p\colon\tilde X\to X$ con $\tilde X$ simplemente conexo. Recuérdese que \emph{simplemente conexo} significa arco-conexo y con grupo fundamental trivial.

Considérese un punto cualquiera $x_0\in X$ y una vecindad regular $U$ de $x_0$, es decir, una vecindad que cumple la definición de espacio cubriente: la preimagen de $U$ es una unión disjunta de abiertos de $\tilde X$ de manera que la restricción de $p$ a cada uno de ellos es un homeomorfismo sobre $U$ (Definiciones~3.1 y~3.2). Sea $\gamma$ un lazo en $x_0$ contenido en la vecindad regular $U$. Este lazo se puede levantar a $\tilde X$: si se levanta el punto inicial a un punto $\tilde x_0$ en alguna de las hojas sobre $U$, el lazo se levanta con la inversa del homeomorfismo (restringido a la vecindad de abajo). Así, $\gamma$ se levanta a un lazo $\tilde\gamma$ dentro de esa hoja, que se ve exactamente idéntico al de abajo, basado en $\tilde x_0$.

Como $\tilde X$ es simplemente conexo, el lazo $\tilde\gamma$ se contrae al lazo constante:
\[
  \tilde\gamma\simeq c_{\tilde x_0}\ \operatorname{rel}\{0,1\}.
\]
Esta misma homotopía, que retrae $\tilde\gamma$ al lazo constante, se puede proyectar con $p$, y da una homotopía de $\gamma$ al lazo constante:
\[
  \gamma\simeq c_{x_0}\ \operatorname{rel}\{0,1\}.
\]

\begin{center}
\begin{tikzpicture}[x=0.022cm,y=-0.022cm,line cap=round]
  \draw[gray] (90,10) -- (400,10) -- (360,110) -- (50,110) -- cycle; \node at (425,30) {$\tilde X$};
  \draw[extremo,thick] (140,40) ellipse (32 and 16); \draw[extremo,thick] (240,35) ellipse (32 and 16);
  \draw[extremo,thick] (130,85) ellipse (32 and 16); \draw[extremo,thick] (250,82) ellipse (32 and 16);
  \draw[thick,azul] (118,87) .. controls (125,72) and (150,72) .. (148,85) .. controls (146,96) and (125,96) .. (118,87);
  \fill[azul] (118,87) circle (2.5); \node[azul,font=\small] at (100,72) {$\tilde x_0$};
  \node[azul,font=\small] at (162,96) {$\tilde\gamma$};
  \draw[->,gray] (220,120) -- (220,160); \node[font=\small] at (233,140) {$p$};
  \draw[gray] (90,170) -- (400,170) -- (360,250) -- (50,250) -- cycle; \node at (425,190) {$X$};
  \draw[extremo,thick] (200,210) ellipse (32 and 16); \node[extremo,font=\small] at (248,200) {$U$};
  \draw[thick,azul] (188,212) .. controls (195,197) and (220,197) .. (218,210) .. controls (216,221) and (195,221) .. (188,212);
  \fill[azul] (188,212) circle (2.5); \node[azul,font=\small] at (175,232) {$x_0$};
  \node[azul,font=\small] at (232,222) {$\gamma$};
\end{tikzpicture}
\end{center}

\noindent Figura: un lazo $\gamma$ dentro de una vecindad regular $U$ se levanta a un lazo $\tilde\gamma$ dentro de una de las hojas (en verde). Como $\tilde X$ es simplemente conexo, $\tilde\gamma\simeq c_{\tilde x_0}$, y al proyectar la homotopía se obtiene $\gamma\simeq c_{x_0}$ en $X$.

Hay que tener cuidado: el dibujo es un poco engañoso, porque pareciera que la homotopía vive dentro de $U$, pero eso no es necesariamente cierto. Para retraer $\tilde\gamma$ al lazo constante podría ser que la homotopía se salga de la hoja, pues nada dice que la hoja sea simplemente conexa (lo que se sabe es que $\tilde X$ lo es); y entonces, al proyectarla, la homotopía también podría salirse de $U$. Lo que se obtiene es que $\gamma$ se retrae al lazo constante en $X$, no necesariamente en $U$.

En resumidas cuentas, se obtiene la siguiente propiedad de $X$: para cada $x\in X$ existe una vecindad $U\subseteq X$ de $x$ tal que todo lazo en $U$ se retrae al lazo constante $c_x$ en $X$ (no necesariamente en $U$, pero sí en $X$). De manera equivalente, para cada $x\in X$ existe una vecindad $U\subseteq X$ de $x$ tal que el homomorfismo inducido por la inclusión
\[
  \pi_1(U,x)\to\pi_1(X,x)
\]
es constante, es decir, todo lazo en $U$ va a dar al elemento trivial. Uno puede convencerse fácilmente de que estas dos propiedades son equivalentes, y se van a usar ambas; sobre todo la de que la inclusión induce un homomorfismo constante.

\begin{defn}[Semilocalmente simplemente conexo]\label{def:slsc}
Un espacio $X$ es \emph{semilocalmente simplemente conexo} si para cada $x\in X$ existe una vecindad $U\subseteq X$ de $x$ tal que el homomorfismo $\pi_1(U,x)\to\pi_1(X,x)$ inducido por la inclusión es constante; equivalentemente, tal que todo lazo en $U$ basado en $x$ se retrae al lazo constante $c_x$ en $X$.
\end{defn}

Lo que se acaba de mostrar es lo siguiente.

\begin{prop}\label{prop:necesaria}
Si $X$ tiene cubriente universal, entonces $X$ es semilocalmente simplemente conexo.
\end{prop}

Así que ser semilocalmente simplemente conexo es una condición necesaria para que $X$ tenga cubriente universal.

\section*{Lo que sigue}

El objetivo de los videos posteriores es mostrar que el regreso también es cierto: si $X$ es arco-conexo, localmente arco-conexo y semilocalmente simplemente conexo, entonces $X$ tiene cubriente universal. La demostración de este teorema es súper constructiva, aunque un poquito tediosa por la notación, porque hay que verificar muchas cosas: el cubriente universal se va a construir completamente a mano. En el siguiente video se habla de la motivación de cómo construirlo.

\end{document}
